\lysh{24347}{2}
\begin{alist}
\item Το πεδίο ορισμού της συνάρτησης $f$ είναι το σύνολο που περιέχει τις τιμές για τις οποίες ορίζεται το κλάσμα, δηλαδή $D_f = \mathbb{R} - \{2\}$.

\item Το κλάσμα της συνάρτησης $f$ απλοποιείται, δηλαδή $f(x) = \frac{x^2-4}{x-2} = \frac{(x-2)(x+2)}{x-2} = x + 2$ για κάθε $x \in \mathbb{R} - \{2\}$.

Άρα, $f(x) = x + 2$.

\item Από το β) ερώτημα έχουμε $\lim\limits_{x\to 2} f(x) = \lim\limits_{x\to 2}(x + 2) = 4$.
\end{alist}
\vspace{6mm}
\noindent\rule{\textwidth}{0.4pt}
\vspace{6mm}

